State some observations about aggregate models and special cases for certain distributios.

Individual risk model:  
\(
S = X_1 + X_2 + \cdots + X_n \rightarrow \mathbb{E}[S] = n\mathbb{E}[X]
\)
\(
\text{Var}(S) = n \text{Var}(X)
\)

Collective risk model:  
\(
S = X_1 + X_2 + \cdots + X_N \rightarrow \mathbb{E}[S] = \mathbb{E}[N]\mathbb{E}[X]
\)
\(
\text{Var}(S) = \mathbb{E}[N] \text{Var}(X) + \text{Var}(N) \mathbb{E}[X]^2
\)

Normal Distribution:  
\(
X_i \overset{iid}{\sim} N(\mu, \sigma^2) \rightarrow S^L = X_1 + X_2 + \cdots + X_n \sim N(n\mu, n\sigma^2)
\)
Exponential Distribution:  
\(
X_i \overset{iid}{\sim} Exp(\theta) \rightarrow S^L = X_1 + X_2 + \cdots + X_n \sim Gamma(\alpha = n, \theta)
\)

Normal approximation:  
\(
S \sim N(\mu = \mathbb{E}[S], \sigma^2 = \text{Var}(S)) \rightarrow \Pr(S \leq k) \approx N\left( \frac{k - \mu}{\sigma} \right)
\)


State the general version, how to determine the distribution of aggregated claims

Let \( S = \sum_{i=1}^N X_i \), where \( X_i \) are i.i.d. and \( N \) is independent of the \( X_i \) and discrete, then
\( \mathbb{P}[S \leq s] = \sum_{n=0}^\infty f_N(n)\mathbb{P}[\sum_{i=1}^N X_i \leq s|N=n] = \sum_{n=0}^\infty f_N(n)F^{\ast n}_X(s)\)
\( \mathbb{P}[S = s] = \sum_{n=0}^\infty f_N(n)f^{\ast n}_X(s)\)
where \( F^{\ast n}(s) \) is the \( n \)-th convolution power, and \( f^{\ast 0}(k) = \delta(k), f^{\ast 1}(k) = f(k) \)

Recall that the distribution of the sum of two r.v. \( X + Y = Z \) is given by the convolution via:
\( f_Z(n) = \sum_{k+m=n} = f_X(k)f_Y(m) = \sum_{k=0}^n f_X(k)f_Y(n-k) \)
\( F_Z(n) = \sum_{k=0}^n = f_X(k)F_Y(n-k) = \sum_{k=0}^n F_X(k)f_Y(n-k) \)
and in the continuous case
\( f_Z(s) = \int_{0}^s f_X(x)f_Y(s - x)dx \)
\( F_Z(s) = \int_{0}^s f_X(x)F_Y(s - x)dx \)
Note that convolution is commutative, so one can exchange \( x,s-x \) in the arguments. 


Define Panjer distributions and give the corresponding recursion scheme

Let \(\mathbf{p} = (p_k)_{k \in \mathbb{N}_0}\) be a sequence of probability weights, i.e.,
\(\forall k \in \mathbb{N}_0: p_k \ge 0\)
\(\sum_{k=0}^\infty p_k = 1\)
\(\mathbf{p}\) is a Panjer distribution if
\(p_0 > 0\)
\(\exists a, b \in \mathbb{R}\) s.t. \(\forall k \in \mathbb{N}: p_k = p_{k-1} \left(a + \frac{b}{k}\right)\)
\(\mathbf{p}\) is degenerate if \(p_0 = 1\).

Some interesting cases are
Binomial
\( a = \frac{-p}{1-p}, b = \frac{p(n+1)}{1-p} \)
Poisson
\( a = 0, b = \lambda \)
NB
\( a = 1-p, b = (1-p)(r-1) \)

\((N, (Y_n)_{n \in \mathbb{N}})\) insurance portfolio from collective model of risk
\(S := \sum_{n=1}^N Y_n\)
\( N \) has nondegenerate Panjer distribution with constants \(a, b \in \mathbb{R}\), \(\mathbf{p} = (\mathbb{P}(N=k))_{k \in \mathbb{N}_0}\)
\(\mathbb{P}_{Y_1}\) supported on \(\mathbb{N}_0\)
\(f_k := \mathbb{P}(S=k)\), \(g_k := \mathbb{P}(Y_1 = k)\), \(k \in \mathbb{N}_0\)
Then
\[
\mathbb{P}(S=r) = f_r = 
\begin{cases}
    \sum_{k=0}^\infty p_k g_0^k & r=0 \\
    \frac{1}{1-ag_0} \sum_{k=1}^r \left(a + \frac{b k}{r}\right) g_r f_{r-k} & r \in \mathbb{N}
\end{cases}
\]

Give examples of compound distributions from the lecture


\(N\) has binomial distribution with parameters \(v, p\) (\(N \sim \text{Bin}(v, p)\)) if
\[
\mathbb{P}(N=n) = \binom{v}{n} p^n (1-p)^{v-n}, \quad n \in \mathcal{A},
\]
where \(\mathcal{A} := \{0, 1, ..., v\}\)

Suppose in the collective model that \(Y_1 \sim G\). 
If \(N \sim \text{Bin}(v, p)\), we say \(S\) has compound binomial distribution (\(S \sim \text{CompBin}(v, p, G)\)).

It's characteristic function is given by \( \phi_S(t) = (p\phi_{Y_1}(t) + 1 - p)^v \).

Let \( v > 0 \) be the deterministic volume and \( \lambda > 0 \) the claims frequency.

\(N\) has Poisson distribution with parameter \(\lambda v > 0\) (\(N \sim \text{Poi}(\lambda v)\)) if
\[
\mathbb{P}(N=n) = e^{-\lambda v} \frac{(\lambda v)^n}{n!}, \quad n \in \mathcal{A},
\]
where \(\mathcal{A} := \mathbb{N}_0\)

Suppose in the collective model that \(Y_1 \sim G\). 
If \(N \sim \text{Poi}(\lambda v)\), we say \(S\) has compound Poisson distribution (\(S \sim \text{CompPoi}(\lambda v, G)\)).


Define Maximum Likelihood Estimation (MLE) and state some special cases

Idea: Given observations \(N_1, ..., N_T\), take parameter under which it is most likely/typical
For \(t = 1, ..., T\) consider
\(
\mathbb{N}_0 \ni k \mapsto p_k^t(\vartheta) := \mathbb{P}^\vartheta(N_t = k)
\)
\(
\Rightarrow \Omega \ni \omega \mapsto p_{N_t(\omega)}^t(\vartheta) \text{ random quantity}
\)
Likelihood function:
\(
L(N_1, ..., N_T, \vartheta) := \prod_{t=1}^T p_{N_t}^t(\vartheta)
\)
Log-likelihood function:
\(
\ell(N_1, ..., N_T, \vartheta) := \log(L(N_1, ..., N_T, \vartheta))
\)
\(
\widehat{\vartheta}(\omega) := \arg \max_{\vartheta \in \Theta} L(N_1, ..., N_T, \vartheta)\) is the MLE.

In case that \( N_1, \dots, N_T \) are independent, then
\( N_t \sim \text{Bin}(v_t, p) \), \( \hat{p} = \frac{\sum_{t=1}^T N_t}{\sum_{t=1}^T v_t} \)
\( N_t \sim \text{Pois}(\lambda v_t) \), \( \hat{\lambda} = \frac{\sum_{t=1}^T N_t}{\sum_{t=1}^T v_t} \)
where \( v_t \) is the volume of policies.

These also coincide with the variance minimizing estimators, which are derived from
\( \hat{\lambda} = \left(\sum_{t = 1}^T \frac{1}{\tau^2_t}\right)^{-1}\sum_{t=1}^T \frac{N_t/v_t}{\tau^2_t} \)
where \( \lambda = \mathbb{E}[\frac{N_t}{v_t}] \) and \( \tau_t^2 = \text{Var}[\frac{N_t}{v_t}] \).
This estimator can't be used directly typically as \( \tau_t^2 \) is not known. 
The corresponding variances for above are
\( \text{Var}[\hat{p}] = \frac{p(1-p)}{\sum_{t=1}^T v_t} \)
\( \text{Var}[\hat{\lambda}] = \frac{\lambda}{\sum_{t=1}^T v_t} \).



Describe the Method of Moments for parameter estimation

The Method of Moments (MoM) estimates parameters \(\theta\) of a probability distribution by equating theoretical moments to sample moments.  
General Procedure:
1. Compute the theoretical moments \( m_k \) in terms of parameters \( \theta \).  
\(
m_k = E[X^k]
\)
2. Compute the sample moments \( \hat{m}_k \) from observed data \( X_1, X_2, ..., X_n \).  
\(
\hat{m}_k = \frac{1}{n} \sum_{i=1}^{n} X_i^k
\)
3. Solve the system of equations: 
\(m_k(\theta_1, \theta_2, \dots, \theta_r) = \hat{m}_k, \quad k = 1, 2, \dots, r\) 
to obtain parameter estimates \( \hat{\theta} \).
## **Examples**  

Example Normal Distribution:
\( X \sim N(\mu, \sigma^2) \)  
Theoretical moments:  
\(
E[X] = \mu, \quad E[X^2] = \mu^2 + \sigma^2
\)
Sample moments:  
\(
\hat{m}_1 = \frac{1}{n} \sum X_i, \quad \hat{m}_2 = \frac{1}{n} \sum X_i^2
\)
Solve for parameters:  
\(
\hat{\mu} = \hat{m}_1
\)
\(
\hat{\sigma}^2 = \hat{m}_2 - \hat{m}_1^2
\)
Example Gamma Distribution:
\( X \sim \text{Gamma}(\alpha, \beta) \)  
Theoretical moments:  
\(
E[X] = \frac{\alpha}{\beta}, \quad E[X^2] = \frac{\alpha(\alpha+1)}{\beta^2}
\)
Sample moments:  
\(
\hat{m}_1 = \frac{1}{n} \sum X_i, \quad \hat{m}_2 = \frac{1}{n} \sum X_i^2
\)
Solve:  
\(
\hat{\beta} = \frac{\hat{m}_1}{\hat{m}_2 - \hat{m}_1^2}, \quad \hat{\alpha} = \hat{\beta} \hat{m}_1
\)

Key Properties
Consistency: As \( n \to \infty \), \( \hat{\theta} \to \theta \) if moments exist and are unique.  
Bias: MoM estimators may be biased.  
Efficiency: MoM is generally less efficient than MLE but simpler.  


Describe the procedure of Layering in insurance

Recall layering from discussion of compound Poisson distribution
\(S = \sum_{i=1}^N Y_i \sim \text{CompPoi}(\lambda v, G), G(M) \in (0, 1) \text{ for } M > 0\)
Then
\[
S_1 := \sum_{i=1}^N Y_i \mathbf{1}_{\{Y_i \le M\}} \text{ and } S_2 := \sum_{i=1}^N Y_i \mathbf{1}_{\{Y_i > M\}}
\]
independent CompPoi-distributed random variables
Importance of layering in insurance: Reinsurance available for different layers.
Fitting claim size data with a single distribution from above potentially not
satisfactory.

Key idea: Fit small and large claims separately
Step 1: Choose \(M > 0\) large s.t. many observed claims \(Y_i\) satisfy \(Y_i \le M\)
Step 2: Fit truncated (i.e. conditioned to \((0, M]\)) distribution from the classes introduced above to data or use empirical distribution
Step 3: Estimate \(G(M)\), e.g. \(G(M) = \frac{\# \text{ claims } \le M}{\# \text{ all claims}}\)
Step 4: Fit Pareto\((M, \alpha)\)-distribution to claims \(Y_i\) with \(Y_i > M\)
Comments on Step 4:

Pareto distribution has heaviest tail
Pareto\((\theta, \alpha)((-\infty, \theta]) = 0\)
Only need to estimate one parameter \(\alpha\) to describe tail

Recall that the Pareto distribution is given by:
\( f_X(x) = \frac{\alpha M^\alpha}{x^{\alpha+1}} \) for \( x \geq M \) otherwise it's \( 0 \)
\( F_X(x) = 1 - (\frac{M}{x})^\alpha \) for \( x \geq M \) otherwise it's \( 0 \)
with
\( \mathbb{E}[X] = \frac{\alpha M}{\alpha - 1} \quad \text{Var}[X] = \frac{\alpha M^2}{(\alpha - 1)^2(\alpha - 2}\) 
for \( \alpha > 1, \alpha > 2 \) respectively otherwise \( \infty \).
